\subsection{A preliminary Euler-residue child in the case $q=5$}

We first eliminate the preliminary Euler-residue child in the branch
\(\min S=5\), namely the branch in which the Euler prime is congruent to
\(2\) modulo \(3\).
This does not yet close the full \(q=5\) branch; the full closure is obtained
after the three \(5\)-input children are closed.

\begin{definition}[The bad Euler branch in the case $\min S=5$]
Let $\mathcal B_{5,\pi\equiv 2(3)}$ denote the branch with the following data:
\(\tau(\mathcal B_{5,\pi\equiv 2(3)})=\mathsf{small}(5)\), so that
\(\min S=5\).  The branch also records that the Euler prime \(\pi\) satisfies
\(\pi\equiv2\pmod3\), and that its exponent \(\alpha\) satisfies the Euler
schema \(\alpha\equiv1\pmod4\).
\end{definition}

\begin{proposition}[Closed $q=5$ branch: Euler prime $2$ modulo $3$]
\label{prop:q5-euler-2mod3-empty}
The branch \(\mathcal B_{5,\pi\equiv 2(3)}\) has no realization by an odd
perfect number.  Equivalently, if \(\min S=5\), then the Euler prime must
satisfy \(\pi\equiv1\pmod3\).
\end{proposition}

\begin{proof}
Assume for contradiction that \(N\) realizes
\(\mathcal B_{5,\pi\equiv 2(3)}\).  Since \(\min S=5\), the prime \(3\) is
not in the support, so lower-prime avoidance with \(\ell=3\) applies to every
prime power in the support.  In particular,
\(\operatorname{Avoid}_3(\pi,\alpha)\) must hold.  But the branch assumes
\(\pi\equiv2\pmod3\), hence \(\ord_3(\pi)=2\).  Since \(\pi\) is the Euler
prime, its exponent satisfies \(\alpha\equiv1\pmod4\).  Thus
\(\alpha+1\equiv2\pmod4\), and in particular \(2\mid\alpha+1\).  By the
cyclotomic input condition, this gives \(3\mid\sigma(\pi^\alpha)\), equivalently
\(v_3(\sigma(\pi^\alpha))>0\).  This contradicts lower-prime avoidance, which
requires \(v_3(\sigma(\pi^\alpha))=0\).
Therefore no odd perfect number realizes the branch
$\mathcal B_{5,\pi\equiv 2(3)}$.
\end{proof}

\begin{corollary}[The minimal prime $5$ is non-Euler]
\label{cor:q5-minimal-prime-non-euler}
If \(N\) is an odd perfect number with \(\min S=5\), then \(5\) is not the
Euler prime.  Consequently \(e_5\equiv0\pmod2\).
\end{corollary}

\begin{proof}
The prime \(5\) satisfies \(5\equiv2\pmod3\).
If $5$ were the Euler prime, then the Euler prime would lie in the refuted branch
of Proposition~\ref{prop:q5-euler-2mod3-empty}.  This is impossible.  Hence $5$
is non-Euler, and by Euler's form its exponent is even:
\(e_5\equiv0\pmod2\).
\end{proof}

\begin{corollary}[First coverage split in the $q=5$ branch]
\label{cor:q5-first-coverage-split}
In the branch \(\mathsf{small}(5)\), the split \(\pi\equiv1\pmod3\) or
\(\pi\equiv2\pmod3\) is complete.  The second child is refuted by
Proposition~\ref{prop:q5-euler-2mod3-empty}.  Hence every surviving
\(q=5\) branch satisfies \(\pi\equiv1\pmod3\).
\end{corollary}

\begin{proof}
Since \(\pi\ne3\) and \(\pi\) is an odd prime, exactly one of
\(\pi\equiv1\pmod3\) and \(\pi\equiv2\pmod3\) holds.
Proposition~\ref{prop:q5-euler-2mod3-empty} refutes the second case.
\end{proof}

\begin{lemma}[Surviving exponent schemas in the branch $\min S=5$]
\label{lem:q5-surviving-exponent-schemas}
Assume that \(N\) is an odd perfect number with \(\min S=5\).  Let \(\pi\) be
the Euler prime, with Euler exponent \(\alpha\).

Then the following constraints hold.

\begin{enumerate}[label=(\roman*),leftmargin=2.2em]

\item The Euler prime satisfies \(\pi\equiv1\pmod3\),
      \(\alpha\equiv1\pmod4\), and \(3\nmid\alpha+1\).

\item The minimal prime \(5\) is non-Euler.  Hence
      \(e_5\equiv0\pmod2\) and \(e_5\ge2\).
      Equivalently, \(e_5+1\) is odd.

\item Let \(p\in S\), \(p\ne5\), be non-Euler.  Then
      \(e_p\equiv0\pmod2\).
      The lower-prime avoidance condition for \(\ell=3\) is
      \(3\nmid e_p+1\) when \(p\equiv1\pmod3\), and is automatic when
      \(p\equiv2\pmod3\).

\item A non-Euler prime \(p\ne5\) is a \(5\)-input prime if and only if
      \(p\equiv1\pmod5\) and \(5\mid e_p+1\).
      In that case \(v_5(\sigma(p^{e_p}))=v_5(e_p+1)\).

\item The Euler prime \(\pi\) is a \(5\)-input prime precisely when either
      \(\pi\equiv1\pmod5\) and \(5\mid\alpha+1\), or
      \(\pi\equiv-1\pmod5\).
      In the first case \(v_5(\sigma(\pi^\alpha))=v_5(\alpha+1)\).
      In the second case
      \[
        v_5(\sigma(\pi^\alpha))
        =
        v_5(\pi^2-1)+v_5\!\left(\frac{\alpha+1}{2}\right).
      \]

\item If \(\pi\equiv2\pmod5\) or \(\pi\equiv3\pmod5\), then
      \(\ord_5(\pi)=4\).  Since \(\alpha+1\equiv2\pmod4\), one has
      \(4\nmid\alpha+1\), so such an Euler prime is \(5\)-neutral.
\end{enumerate}
\end{lemma}

\begin{proof}
The first assertion follows from
Proposition~\ref{prop:q5-euler-2mod3-empty}.  Since $\pi\equiv1\pmod3$, one has
\(\ord_3(\pi)=1\).  Lower-prime avoidance for \(\ell=3\) therefore gives
\(3\nmid\alpha+1\).  The congruence \(\alpha\equiv1\pmod4\) is Euler's parity
condition.

The second assertion is Corollary~\ref{cor:q5-minimal-prime-non-euler}.

Now let $p\ne5$ be non-Euler.  Then Euler's form gives
\(e_p\equiv0\pmod2\).
If $p\equiv1\pmod3$, then $\ord_3(p)=1$, so lower-prime avoidance gives
\(3\nmid e_p+1\).
If $p\equiv2\pmod3$, then $\ord_3(p)=2$.  Since $e_p$ is even,
\(e_p+1\) is odd, so \(2\nmid e_p+1\).
Thus the avoidance condition is automatic in this case.

For the $5$-input classification, apply the cyclotomic input condition with
$q=5$.  If $p$ is non-Euler, then $e_p+1$ is odd.  The possible exact orders
modulo $5$ are
\[
  1,\ 2,\ 4.
\]
The even orders $2$ and $4$ cannot divide the odd integer $e_p+1$.  Thus a
non-Euler \(5\)-input can only have \(\ord_5(p)=1\), i.e.
\(p\equiv1\pmod5\).
In the order-one case, the corrected cyclotomic input condition says that input
occurs precisely when \(5\mid e_p+1\), and the contribution is
\(v_5(\sigma(p^{e_p}))=v_5(e_p+1)\).

For the Euler prime, put \(n_\pi:=\alpha+1\).  Since
\(\alpha\equiv1\pmod4\), we have \(n_\pi\equiv2\pmod4\).
If $\ord_5(\pi)=1$, then the order-one case gives input precisely when
\(5\mid n_\pi\).
If $\ord_5(\pi)=2$, then $2\mid n_\pi$, so the cyclotomic input condition gives
a $5$-adic contribution
\[
  v_5(\pi^2-1)+v_5(n_\pi/2).
\]
If $\ord_5(\pi)=4$, then $4\nmid n_\pi$, so there is no $5$-adic contribution.
This proves the classification.
\end{proof}

\begin{corollary}[Closed $q=5$ leaf: no $5$-input prime]
\label{cor:q5-no-input-leaf-empty}
Assume
\[
  \min S=5.
\]
The child branch in which every prime
\[
  p\in S,\qquad p\ne5,
\]
is $5$-neutral has no realization by an odd perfect number.
\end{corollary}

\begin{proof}
Since $5\in S$, one has
\[
  e_5>0.
\]
The $5$-adic valuation balance is
\[
  e_5
  =
  \sum_{\substack{p\in S\\p\ne5}}
  v_5(\sigma(p^{e_p})).
\]
In the stated child branch every prime $p\ne5$ is $5$-neutral, so every summand
on the right-hand side is zero.  Hence
\[
  e_5=0,
\]
contradicting
\[
  5\in S.
\]
Therefore the branch is empty.
\end{proof}

\subsubsection{The first necessary $5$-input split}

We now turn the $5$-adic balance in the branch
\[
  \min S=5
\]
into a first explicit coverage split.

\begin{definition}[$5$-input witness in the branch $\min S=5$]
Assume
\[
  \min S=5.
\]
A prime
\[
  p\in S,\qquad p\ne5,
\]
is called a $5$-input witness if
\[
  v_5(\sigma(p^{e_p}))>0.
\]
Equivalently, $p$ is a prime power in the support, other than $5^{e_5}$ itself,
that contributes to the $5$-adic balance
\[
  e_5
  =
  \sum_{\substack{p\in S\\p\ne5}}
  v_5(\sigma(p^{e_p})).
\]
\end{definition}

\begin{proposition}[Existence of a $5$-input witness]
\label{prop:q5-input-witness-exists}
Assume that $N$ is an odd perfect number with
\[
  \min S=5.
\]
Then there exists at least one $5$-input witness
\[
  p\in S,\qquad p\ne5.
\]
\end{proposition}

\begin{proof}
Since
\[
  5\in S,
\]
one has
\[
  e_5>0.
\]
The $5$-adic valuation balance gives
\[
  e_5
  =
  \sum_{\substack{p\in S\\p\ne5}}
  v_5(\sigma(p^{e_p})).
\]
The right-hand side is a sum of nonnegative integers.  Since the left-hand side
is positive, at least one summand is positive.  Hence there exists
\[
  p\in S,\qquad p\ne5,
\]
with
\[
  v_5(\sigma(p^{e_p}))>0.
\]
\end{proof}

\begin{lemma}[Classification of first $5$-input witnesses]
\label{lem:q5-first-input-classification}
Assume
\[
  \min S=5.
\]
Let
\[
  p\in S,\qquad p\ne5,
\]
be a $5$-input witness.

Then exactly one of the following role cases applies.

\begin{enumerate}[label=(\roman*),leftmargin=2.2em]

\item \textbf{Non-Euler witness.}
      The prime $p$ is non-Euler.  Then
      \[
        e_p\equiv0\pmod2,
      \]
      and the $5$-input condition is equivalent to
      \[
        p\equiv1\pmod5,
        \qquad
        5\mid e_p+1.
      \]
      In this case
      \[
        v_5(\sigma(p^{e_p}))=v_5(e_p+1).
      \]

\item \textbf{Euler witness of order one.}
      The prime $p$ is the Euler prime $\pi$, and
      \[
        \pi\equiv1\pmod5,
        \qquad
        5\mid \alpha+1.
      \]
      In this case
      \[
        v_5(\sigma(\pi^\alpha))=v_5(\alpha+1).
      \]

\item \textbf{Euler witness of order two.}
      The prime $p$ is the Euler prime $\pi$, and
      \[
        \pi\equiv-1\pmod5.
      \]
      In this case
      \[
        \ord_5(\pi)=2,
      \]
      and
      \[
        v_5(\sigma(\pi^\alpha))
        =
        v_5(\pi^2-1)
        +
        v_5\!\left(\frac{\alpha+1}{2}\right).
      \]
\end{enumerate}

Moreover, the Euler witness cannot have
\[
  \pi\equiv2\pmod5
  \qquad\text{or}\qquad
  \pi\equiv3\pmod5,
\]
because then
\[
  \ord_5(\pi)=4
\]
and
\[
  4\nmid \alpha+1.
\]
\end{lemma}

\begin{proof}
Let
\[
  n_p:=e_p+1.
\]

First suppose that $p$ is non-Euler.  Then Euler's form gives
\[
  e_p\equiv0\pmod2,
\]
so
\[
  n_p=e_p+1
\]
is odd.  Since $p$ is a $5$-input witness, the cyclotomic input condition gives
\[
  \ord_5(p)\mid n_p
\]
if $\ord_5(p)>1$, and gives
\[
  5\mid n_p
\]
if $\ord_5(p)=1$.

The possible orders modulo $5$ are
\[
  1,\ 2,\ 4.
\]
The orders $2$ and $4$ are even and therefore cannot divide the odd integer
$n_p$.  Hence the only possible non-Euler input order is
\[
  \ord_5(p)=1,
\]
which is equivalent to
\[
  p\equiv1\pmod5.
\]
In the order-one case the corrected cyclotomic input condition gives input
precisely when
\[
  5\mid e_p+1,
\]
and the valuation is
\[
  v_5(\sigma(p^{e_p}))=v_5(e_p+1).
\]

Now suppose that $p=\pi$ is the Euler prime, with exponent $\alpha$.  Put
\[
  n_\pi:=\alpha+1.
\]
Since
\[
  \alpha\equiv1\pmod4,
\]
we have
\[
  n_\pi\equiv2\pmod4.
\]

If
\[
  \ord_5(\pi)=1,
\]
then
\[
  \pi\equiv1\pmod5,
\]
and the order-one case gives input precisely when
\[
  5\mid \alpha+1.
\]
The valuation is then
\[
  v_5(\sigma(\pi^\alpha))=v_5(\alpha+1).
\]

If
\[
  \ord_5(\pi)=2,
\]
then
\[
  \pi\equiv-1\pmod5.
\]
Since
\[
  2\mid n_\pi,
\]
the cyclotomic input condition gives a $5$-adic contribution, and the valuation
formula is
\[
  v_5(\sigma(\pi^\alpha))
  =
  v_5(\pi^2-1)
  +
  v_5\!\left(\frac{\alpha+1}{2}\right).
\]

Finally, if
\[
  \ord_5(\pi)=4,
\]
then
\[
  \pi\equiv2\pmod5
  \quad\text{or}\quad
  \pi\equiv3\pmod5.
\]
But
\[
  n_\pi\equiv2\pmod4,
\]
so
\[
  4\nmid n_\pi.
\]
Thus no $5$-adic contribution occurs in these two residue classes.
\end{proof}

\begin{corollary}[First necessary $5$-input coverage split]
\label{cor:q5-first-input-coverage-split}
In the branch
\[
  \min S=5,
\]
the following split is complete:
\[
  \text{there exists a non-Euler $5$-input witness}
\]
or
\[
  \text{the Euler prime is a $5$-input witness}.
\]
More explicitly, every odd perfect number with \(\min S=5\) realizes at least
one of the following child branches: \(\mathbf{Q5N}\), where there is a
non-Euler \(p\in S\) with \(p\ne5,\ p\ne\pi,\ p\equiv1\pmod5\), and
\(5\mid e_p+1\); \(\mathbf{Q5E1}\), where \(\pi\equiv1\pmod5\) and
\(5\mid\alpha+1\); or \(\mathbf{Q5E2}\), where \(\pi\equiv-1\pmod5\).  The
complementary child branch, in which every prime \(p\in S\), \(p\ne5\), is
\(5\)-neutral, is empty.
\end{corollary}

\begin{proof}
By Proposition~\ref{prop:q5-input-witness-exists}, at least one $5$-input
witness exists.

If some $5$-input witness is non-Euler, then
Lemma~\ref{lem:q5-first-input-classification} gives the child branch
\(\mathbf{Q5N}\).

If no non-Euler witness exists, then some $5$-input witness must be the Euler
prime $\pi$.  Lemma~\ref{lem:q5-first-input-classification} shows that the Euler
witness is either of order one modulo \(5\), giving \(\mathbf{Q5E1}\), or of
order two modulo \(5\), giving \(\mathbf{Q5E2}\).

The branch in which all primes $p\ne5$ are $5$-neutral is exactly the branch
refuted by Corollary~\ref{cor:q5-no-input-leaf-empty}.
\end{proof}

\begin{remark}[Non-disjointness of the first $5$-input split]
The split in Corollary~\ref{cor:q5-first-input-coverage-split} is a coverage
split, not necessarily a disjoint split.  For example, a hypothetical branch
could contain both a non-Euler $5$-input witness and an Euler $5$-input witness.
For the finite search, this is harmless: coverage requires that every
hypothetical odd perfect number realizes at least one child branch, not that it
realizes exactly one.
\end{remark}

\begin{remark}[Certificate data for the first $5$-input split]
A certificate for the first $5$-input split records one of the following witness
types: \(\mathtt{Q5N}\), \(\mathtt{Q5E1}\), or \(\mathtt{Q5E2}\).
For a \(\mathtt{Q5N}\) record the verifier checks
\[
  p\ne5,\quad p\ne\pi,\quad p\equiv1\pmod5,\quad 5\mid e_p+1,
  \quad e_p\equiv0\pmod2.
\]
For a \(\mathtt{Q5E1}\) record it checks
\[
  \pi\equiv1\pmod5,\quad 5\mid\alpha+1,\quad \alpha\equiv1\pmod4.
\]
For a \(\mathtt{Q5E2}\) record it checks
\[
  \pi\equiv-1\pmod5,\quad \alpha\equiv1\pmod4.
\]
In each case the cyclotomic input condition proves
\[
  v_5(\sigma(p^{e_p}))>0
\]
for the recorded witness.
\end{remark}

\begin{proposition}[Completeness of the first \(q=5\) coverage split]
\label{prop:q5-first-input-three-child-coverage}
Assume that \(N\) is an odd perfect number with \(\min\{p:p\mid N\}=5\).
Then \(N\) realizes at least one of the following three branch families:
\(\mathbf{Q5E2}\), \(\mathbf{Q5E1}\), and \(\mathbf{Q5N}\).
More explicitly, at least one of the following alternatives holds.

\[
\begin{array}{ll}
\mathbf{Q5E2}: & \pi\equiv-1\pmod5,\\[2pt]
\mathbf{Q5E1}: & \pi\equiv1\pmod5,\quad 5\mid\alpha+1,\\[2pt]
\mathbf{Q5N}: &
\text{there is a non-Euler }p\in S,\ p\ne5,\ p\ne\pi,\ 
p\equiv1\pmod5,\ 5\mid e_p+1.
\end{array}
\]

Thus the three children \(\mathbf{Q5E2}\), \(\mathbf{Q5E1}\), and
\(\mathbf{Q5N}\) form a complete coverage split of the \(q=5\) minimal-prime
branch.  The split is a coverage split, not necessarily a disjoint split.
\end{proposition}

\begin{proof}
Assume \(\min\{p:p\mid N\}=5\).  Then \(5\in S\), so \(e_5>0\).  The
\(5\)-adic valuation balance for \(\sigma(N)=2N\) is
\[
  e_5
  =
  \sum_{\substack{p\in S\\p\ne5}}
  v_5(\sigma(p^{e_p})).
\]
The right-hand side is a sum of nonnegative integers.  Since \(e_5>0\), at least
one summand is positive.  Hence there exists a prime \(p\in S\), \(p\ne5\),
with \(v_5(\sigma(p^{e_p}))>0\).  We call such a prime a \(5\)-input witness.

We now distinguish whether the witness is the Euler prime or not.

First suppose the witness is not the Euler prime.  Then \(p\ne\pi\), and Euler's
form gives \(e_p\equiv0\pmod2\).  Thus \(e_p+1\) is odd.  Since
\(v_5(\sigma(p^{e_p}))>0\),
the cyclotomic input condition gives either
\(\operatorname{ord}_5(p)\mid e_p+1\) if \(\operatorname{ord}_5(p)>1\), or
\(5\mid e_p+1\) if \(\operatorname{ord}_5(p)=1\).

The possible exact orders modulo \(5\) are
\[
  1,\ 2,\ 4.
\]
The orders \(2\) and \(4\) cannot divide the odd integer \(e_p+1\).  Therefore
\(\operatorname{ord}_5(p)=1\), equivalently \(p\equiv1\pmod5\).  The order-one
input condition then gives \(5\mid e_p+1\).  This is precisely the branch
\(\mathbf{Q5N}\).

It remains to treat the case in which every \(5\)-input witness is the Euler
prime \(\pi\).  Let \(\alpha=e_\pi\) be its exponent.  Euler's form gives
\(\alpha\equiv1\pmod4\), hence \(\alpha+1\equiv2\pmod4\).  Since \(\pi\) is a
\(5\)-input witness, \(v_5(\sigma(\pi^\alpha))>0\).  Let
\(d:=\operatorname{ord}_5(\pi)\).
The possible exact orders modulo \(5\) are again
\[
  d\in\{1,2,4\}.
\]

If \(d=1\), then \(\pi\equiv1\pmod5\), and the order-one cyclotomic input
condition gives \(5\mid\alpha+1\); this is exactly \(\mathbf{Q5E1}\).  If
\(d=2\), then \(\pi\equiv-1\pmod5\), which is exactly \(\mathbf{Q5E2}\).  Finally,
if \(d=4\), then a \(5\)-adic input from \(\sigma(\pi^\alpha)\) would require
\(4\mid\alpha+1\), impossible because \(\alpha+1\equiv2\pmod4\).  Hence the
Euler witness cannot have order \(4\) modulo \(5\).

Thus every possible \(5\)-input witness leads to one of \(\mathbf{Q5N}\),
\(\mathbf{Q5E1}\), or \(\mathbf{Q5E2}\).  Therefore the three displayed
children cover all realizations of the branch \(\min\{p:p\mid N\}=5\).

The split need not be disjoint.  A hypothetical branch could contain both a
non-Euler \(5\)-input witness and an Euler \(5\)-input witness.  For branch
coverage this is harmless: completeness requires that every realization lies in
at least one child, not in exactly one child.
\end{proof}

\begin{corollary}[Reduction of the \(q=5\) branch to the three certified children]
\label{cor:q5-reduction-to-three-certified-children}
To exhaust the \(\CertRelease\) \(q=5\) minimal-prime branch it is enough to close
the three child inventories \(\mathbf{Q5E2}\), \(\mathbf{Q5E1}\), and
\(\mathbf{Q5N}\).
In particular, if each of these three child inventories has no remaining
reduced, unresolved, or delegated terminal obligation, then the \(\CertRelease\)
\(q=5\) branch inventory is exhausted.
\end{corollary}

\begin{proof}
By Proposition~\ref{prop:q5-first-input-three-child-coverage}, every odd perfect
number in the branch
\(\min\{p:p\mid N\}=5\) realizes at least one of \(\mathbf{Q5E2}\),
\(\mathbf{Q5E1}\), and \(\mathbf{Q5N}\).
If all three child inventories are closed, then no realization remains in any
child.  Since the children cover the parent branch, no realization remains in
the \(q=5\) branch.
\end{proof}



\subsubsection{Forced extensions in the first $5$-input branches}

We now attach the first forced cyclotomic obligations to the three child
branches \(\mathbf{Q5N}\), \(\mathbf{Q5E1}\), and \(\mathbf{Q5E2}\) from
Corollary~\ref{cor:q5-first-input-coverage-split}.

\begin{lemma}[Forced extension in the branch $\mathbf{Q5N}$]
\label{lem:q5n-forced-extension}
Assume \(\min S=5\).  Let \(p\in S\), \(p\ne5\), \(p\ne\pi\), be a non-Euler
\(5\)-input witness in the branch \(\mathbf{Q5N}\).  Thus
\(p\equiv1\pmod5\) and \(5\mid e_p+1\).
Then
\[
  \Phi_5(p)=p^4+p^3+p^2+p+1
\]
divides \(\sigma(p^{e_p})\), and \(v_5(\Phi_5(p))=1\).
Consequently
\[
  C_5(p):=\frac{\Phi_5(p)}5
\]
is an integer with \(C_5(p)>1\).

If \(C_5(p)\) has a prime divisor equal to \(3\), then the branch is immediately
refuted by lower-prime avoidance.  Otherwise \(C_5(p)\) has a prime divisor
\(r>5\), and every such divisor is forced into the support, with \(r\in S\).
Moreover, for every prime divisor \(r\mid C_5(p)\) with \(r>5\), one has
\(\ord_r(p)=5\), and hence \(r\equiv1\pmod5\).
\end{lemma}

\begin{proof}
Since \(5\mid e_p+1\), the cyclotomic factor \(\Phi_5(p)\) divides
\(\sigma(p^{e_p})\).  Because \(p\equiv1\pmod5\), LTE gives
\(v_5(p^5-1)=v_5(p-1)+1\).  Since \(\Phi_5(p)=(p^5-1)/(p-1)\), we obtain
\(v_5(\Phi_5(p))=1\).  Thus \(C_5(p)=\Phi_5(p)/5\) is an integer.  Also
\(\Phi_5(p)>5\) for every prime \(p\ge7\), and in the branch \(\min S=5\) we
have \(p\ne5\).  Hence \(C_5(p)>1\).

Let \(r\mid C_5(p)\) be prime.  Then
\(r\mid\Phi_5(p)\mid\sigma(p^{e_p})\mid\sigma(N)=2N\).  Since \(r\) is odd,
this gives \(r\mid N\), so \(r\in S\).  If \(r=3\), this contradicts
\(\min S=5\), because \(3\notin S\).  Hence any surviving branch has a prime
divisor \(r>5\) of \(C_5(p)\), and such an \(r\) is forced into \(S\).

For \(r>5\), we have \(r\nmid5\).  Also \(r\ne p\), because
\(\Phi_5(p)=p^4+p^3+p^2+p+1\equiv1\pmod p\).  Since \(r\mid\Phi_5(p)\), the
exact order of \(p\) modulo \(r\) is \(\ord_r(p)=5\).  Therefore \(5\mid r-1\),
so \(r\equiv1\pmod5\).
\end{proof}

\begin{lemma}[Forced extension in the branch $\mathbf{Q5E1}$]
\label{lem:q5e1-forced-extension}
Assume \(\min S=5\).
Suppose the Euler prime satisfies the branch condition
\[
  \mathbf{Q5E1}:\qquad
  \pi\equiv1\pmod5,
  \qquad
  5\mid\alpha+1.
\]
Then
\[
  \Phi_5(\pi)=\pi^4+\pi^3+\pi^2+\pi+1
\]
divides \(\sigma(\pi^\alpha)\), and \(v_5(\Phi_5(\pi))=1\).
Consequently
\[
  C_5(\pi):=\frac{\Phi_5(\pi)}5
\]
is an integer with \(C_5(\pi)>1\).

If \(C_5(\pi)\) has a prime divisor equal to \(3\), then the branch is
immediately refuted by lower-prime avoidance.  Otherwise \(C_5(\pi)\) has a
prime divisor \(r>5\), and every such divisor is forced into the support, with
\(r\in S\).  Moreover, for every prime divisor \(r\mid C_5(\pi)\) with \(r>5\),
one has \(\ord_r(\pi)=5\), and hence \(r\equiv1\pmod5\).
\end{lemma}

\begin{proof}
The proof is identical to Lemma~\ref{lem:q5n-forced-extension}, with \(p\)
replaced by the Euler prime \(\pi\).  Since \(5\mid\alpha+1\), the factor
\(\Phi_5(\pi)\) divides \(\sigma(\pi^\alpha)\).  Since \(\pi\equiv1\pmod5\),
LTE gives \(v_5(\Phi_5(\pi))=1\).  Thus \(C_5(\pi)=\Phi_5(\pi)/5\) is a
positive integer.

If a prime divisor of \(C_5(\pi)\) is \(3\), the branch is refuted because
\(3\notin S\) when \(\min S=5\).
Otherwise \(C_5(\pi)\) has a prime divisor \(r>5\), and every such divisor is
forced into \(S\).  Since \(r\mid\Phi_5(\pi)\), \(r\nmid5\), and
\(\Phi_5(\pi)\equiv1\pmod\pi\), we have \(r\ne\pi\) and \(\ord_r(\pi)=5\).
Hence \(r\equiv1\pmod5\).
\end{proof}

\begin{lemma}[No pure $\mathbf{Q5E2}$ exceptional branch]
\label{lem:no-pure-q5e2-exception}
There is no prime \(\pi\equiv1\pmod4\) and no integer \(c\ge1\) such that
\[
  \pi+1=5^c.
\]
Consequently, in the branch
\(\mathbf{Q5E2}\), where \(\pi\equiv-1\pmod5\),
the cofactor obtained from \(\pi+1\) after removing its full \(5\)-power is
never equal to \(1\).
\end{lemma}

\begin{proof}
Assume \(\pi+1=5^c\) with \(\pi\equiv1\pmod4\).  Then
\(\pi+1\equiv2\pmod4\), but \(5^c\equiv1\pmod4\) for every \(c\ge1\).  This is
impossible.  Hence no such \(\pi\) and \(c\) exist.
\end{proof}

\begin{lemma}[Forced extension in the branch $\mathbf{Q5E2}$]
\label{lem:q5e2-forced-extension}
Assume \(\min S=5\).
Suppose the Euler prime satisfies the branch condition
\(\mathbf{Q5E2}\), namely \(\pi\equiv-1\pmod5\).
Then
\[
  \Phi_2(\pi)=\pi+1
\]
divides \(\sigma(\pi^\alpha)\), and \(v_5(\pi+1)\ge1\).
Put
\[
  c_\pi:=v_5(\pi+1),
  \qquad
  C_2(\pi):=\frac{\pi+1}{5^{c_\pi}}.
\]
Then \(C_2(\pi)>1\).

If \(C_2(\pi)\) has a prime divisor equal to \(3\), then the branch is
immediately refuted by lower-prime avoidance.  Otherwise \(C_2(\pi)\) has a
prime divisor \(r>5\), and every such divisor is forced into the support, with
\(r\in S\).  Moreover, for every prime divisor \(r\mid C_2(\pi)\) with \(r>5\),
one has \(\pi\equiv-1\pmod r\), so \(\ord_r(\pi)=2\), and therefore
\(r\equiv1\pmod2\).
\end{lemma}

\begin{proof}
Since \(\pi\) is the Euler prime, \(\alpha\equiv1\pmod4\), hence
\(2\mid\alpha+1\).  Therefore \(\Phi_2(\pi)=\pi+1\) divides
\(\sigma(\pi^\alpha)\).  The branch condition \(\pi\equiv-1\pmod5\) gives
\(5\mid\pi+1\), so \(c_\pi=v_5(\pi+1)\ge1\).  If \(C_2(\pi)=1\), then
\(\pi+1=5^{c_\pi}\), contradicting Lemma~\ref{lem:no-pure-q5e2-exception}.
Hence \(C_2(\pi)>1\).

Let \(r\mid C_2(\pi)\) be prime.  Then
\(r\mid\pi+1\mid\sigma(\pi^\alpha)\mid\sigma(N)=2N\).  Since \(r\) is odd,
this gives \(r\mid N\), so \(r\in S\).  If \(r=3\), this contradicts
\(\min S=5\).  Thus every surviving branch has a prime divisor \(r>5\) of
\(C_2(\pi)\), and such an \(r\) is forced into \(S\).

For every such \(r\), one has \(\pi\equiv-1\pmod r\).  Since \(r>5\), in
particular \(r\ne2\), and therefore \(\ord_r(\pi)=2\).
\end{proof}

\begin{corollary}[Certified first-step reduction after the $5$-input split]
\label{cor:q5-first-input-forced-reduction}
In the branch
\[
  \min S=5,
\]
each child branch from the first \(5\)-input split has the following certified
outcome.

\begin{enumerate}[label=(\roman*),leftmargin=2.2em]

\item In \(\mathbf{Q5N}\), the witness \(p\) forces the cofactor
      \[
        C_5(p)=\frac{\Phi_5(p)}5.
      \]
      If \(3\mid C_5(p)\), the branch is refuted by lower-prime avoidance.
      Otherwise a prime divisor
      \[
        r>5,\qquad r\equiv1\pmod5,
      \]
      is forced into \(S\).

\item In \(\mathbf{Q5E1}\), the Euler prime \(\pi\) forces the cofactor
      \[
        C_5(\pi)=\frac{\Phi_5(\pi)}5.
      \]
      If \(3\mid C_5(\pi)\), the branch is refuted by lower-prime avoidance.
      Otherwise a prime divisor
      \[
        r>5,\qquad r\equiv1\pmod5,
      \]
      is forced into \(S\).

\item In \(\mathbf{Q5E2}\), the Euler prime \(\pi\) forces the cofactor
      \[
        C_2(\pi)=\frac{\pi+1}{5^{v_5(\pi+1)}}.
      \]
      The pure case \(C_2(\pi)=1\) is impossible by
      Lemma~\ref{lem:no-pure-q5e2-exception}.  If \(3\mid C_2(\pi)\), the branch
      is refuted by lower-prime avoidance.  Otherwise a prime divisor
      \[
        r>5
      \]
      is forced into \(S\), and it satisfies
      \[
        \ord_r(\pi)=2.
      \]
\end{enumerate}
\end{corollary}

\begin{proof}
The three cases are exactly Lemmas~\ref{lem:q5n-forced-extension},
\ref{lem:q5e1-forced-extension}, and \ref{lem:q5e2-forced-extension}.
\end{proof}

\begin{remark}[Certificate data for the forced $5$-input extensions]
A certificate for the forced extension after the first \(5\)-input split records
one of the following types:
\[
  \mathtt{Q5N\_Phi5},
  \qquad
  \mathtt{Q5E1\_Phi5},
  \qquad
  \mathtt{Q5E2\_Phi2}.
\]
For the two \(\Phi_5\)-records the verifier checks
\[
  \Phi_5(x)=x^4+x^3+x^2+x+1,
  \qquad
  v_5(\Phi_5(x))=1,
\]
where \(x=p\) in \(\mathbf{Q5N}\) and \(x=\pi\) in \(\mathbf{Q5E1}\).
It then factors
\[
  \Phi_5(x)/5.
\]
A divisor \(3\) gives a lower-prime avoidance refutation, while any divisor
\(r>5\) is a forced support prime with
\[
  \ord_r(x)=5.
\]

For the \(\mathbf{Q5E2}\)-record the verifier checks
\[
  \Phi_2(\pi)=\pi+1,
  \qquad
  c_\pi=v_5(\pi+1),
\]
and factors
\[
  (\pi+1)/5^{c_\pi}.
\]
The case in which this cofactor is \(1\) is disallowed by
Lemma~\ref{lem:no-pure-q5e2-exception}.  A divisor \(3\) gives a lower-prime
avoidance refutation, while any divisor \(r>5\) is a forced support prime with
\[
  \ord_r(\pi)=2.
\]
\end{remark}

\subsubsection{Second-level coverage splits after the forced $5$-input extensions}

We now refine the forced extensions obtained from the first \(5\)-input split.
The goal is to turn each forced cofactor into a certified second-level branch:
either a lower-prime avoidance contradiction occurs, or a new support prime is
forced.

\begin{definition}[Second-level cofactors in the \(q=5\) branch]
In the branches \(\mathbf{Q5N}\) and \(\mathbf{Q5E1}\), put
\[
  C_5(x):=\frac{\Phi_5(x)}5
  =
  \frac{x^4+x^3+x^2+x+1}{5},
\]
where
\[
  x=p
  \quad\text{in }\mathbf{Q5N},
  \qquad
  x=\pi
  \quad\text{in }\mathbf{Q5E1}.
\]

In the branch \(\mathbf{Q5E2}\), put
\[
  C_2(\pi):=
  \frac{\pi+1}{5^{v_5(\pi+1)}}.
\]
\end{definition}

\begin{proposition}[Second-level split in the branches \(\mathbf{Q5N}\) and \(\mathbf{Q5E1}\)]
\label{prop:q5-phi5-second-level-split}
Assume
\[
  \min S=5.
\]
Let \(x\) denote either the non-Euler witness \(p\) in \(\mathbf{Q5N}\), or the
Euler prime \(\pi\) in \(\mathbf{Q5E1}\).  Thus
\[
  x\equiv1\pmod5,
  \qquad
  5\mid e_x+1,
\]
where \(e_x=e_p\) in \(\mathbf{Q5N}\), and \(e_x=\alpha\) in
\(\mathbf{Q5E1}\).

Then the split
\[
  3\mid C_5(x)
  \qquad\text{or}\qquad
  3\nmid C_5(x)
\]
is complete.

If
\[
  3\mid C_5(x),
\]
then the branch is refuted by lower-prime avoidance.

If
\[
  3\nmid C_5(x),
\]
then there exists a prime
\[
  r>5
\]
such that
\[
  r\mid C_5(x),
  \qquad
  r\in S,
\]
and
\[
  \ord_r(x)=5.
\]
Consequently
\[
  r\equiv1\pmod5.
\]
\end{proposition}

\begin{proof}
By Lemmas~\ref{lem:q5n-forced-extension} and
\ref{lem:q5e1-forced-extension}, the cofactor
\[
  C_5(x)=\frac{\Phi_5(x)}5
\]
is an integer and satisfies
\[
  C_5(x)>1.
\]
Therefore \(C_5(x)\) has at least one prime divisor.

The split
\[
  3\mid C_5(x)
  \qquad\text{or}\qquad
  3\nmid C_5(x)
\]
is exhaustive.

If
\[
  3\mid C_5(x),
\]
then
\[
  3\mid \Phi_5(x)\mid \sigma(x^{e_x}).
\]
Hence
\[
  3\mid \sigma(N)=2N.
\]
Since \(3\) is odd, this forces
\[
  3\mid N,
\]
contradicting
\[
  \min S=5.
\]
Thus this child branch is refuted by lower-prime avoidance.

Now suppose
\[
  3\nmid C_5(x).
\]
Since \(C_5(x)>1\), choose a prime divisor
\[
  r\mid C_5(x).
\]
The divisor \(r\) is not \(3\) by assumption.  It is also not \(5\), because the
full \(5\)-factor of \(\Phi_5(x)\) has been removed:
\[
  v_5(\Phi_5(x))=1.
\]
Therefore
\[
  r>5.
\]

Since
\[
  r\mid C_5(x)\mid \Phi_5(x)\mid \sigma(x^{e_x})\mid \sigma(N)=2N,
\]
and \(r\) is odd, we get
\[
  r\mid N.
\]
Thus
\[
  r\in S.
\]

Moreover,
\[
  r\ne x,
\]
because
\[
  \Phi_5(x)=x^4+x^3+x^2+x+1\equiv1\pmod x.
\]
Since
\[
  r\mid\Phi_5(x)
\]
and
\[
  r\nmid5,
\]
the exact order of \(x\) modulo \(r\) is
\[
  \ord_r(x)=5.
\]
Hence
\[
  5\mid r-1,
\]
so
\[
  r\equiv1\pmod5.
\]
\end{proof}

\begin{proposition}[Second-level split in the branch \(\mathbf{Q5E2}\)]
\label{prop:q5e2-second-level-split}
Assume
\[
  \min S=5.
\]
Suppose the Euler prime satisfies
\[
  \mathbf{Q5E2}:\qquad \pi\equiv-1\pmod5.
\]
Put
\[
  C_2(\pi):=\frac{\pi+1}{5^{v_5(\pi+1)}}.
\]
Then the split
\[
  3\mid C_2(\pi)
  \qquad\text{or}\qquad
  3\nmid C_2(\pi)
\]
is complete.

If
\[
  3\mid C_2(\pi),
\]
then the branch is refuted by lower-prime avoidance.

If
\[
  3\nmid C_2(\pi),
\]
then there exists a prime
\[
  r>5
\]
such that
\[
  r\mid C_2(\pi),
  \qquad
  r\in S,
\]
and
\[
  \ord_r(\pi)=2.
\]
\end{proposition}

\begin{proof}
By Lemma~\ref{lem:q5e2-forced-extension}, the cofactor
\[
  C_2(\pi)=\frac{\pi+1}{5^{v_5(\pi+1)}}
\]
is an integer satisfying
\[
  C_2(\pi)>1.
\]
Thus \(C_2(\pi)\) has at least one prime divisor.

The split
\[
  3\mid C_2(\pi)
  \qquad\text{or}\qquad
  3\nmid C_2(\pi)
\]
is exhaustive.

If
\[
  3\mid C_2(\pi),
\]
then
\[
  3\mid \pi+1=\Phi_2(\pi)\mid\sigma(\pi^\alpha).
\]
Hence
\[
  3\mid\sigma(N)=2N.
\]
Since \(3\) is odd, this forces
\[
  3\mid N,
\]
contradicting
\[
  \min S=5.
\]
Thus this child branch is refuted by lower-prime avoidance.

Now suppose
\[
  3\nmid C_2(\pi).
\]
Since \(C_2(\pi)>1\), choose a prime divisor
\[
  r\mid C_2(\pi).
\]
The divisor \(r\) is not \(3\) by assumption, and it is not \(5\) by the
definition of \(C_2(\pi)\).  Therefore
\[
  r>5.
\]

Since
\[
  r\mid C_2(\pi)\mid \pi+1\mid\sigma(\pi^\alpha)\mid\sigma(N)=2N,
\]
and \(r\) is odd, we get
\[
  r\mid N.
\]
Thus
\[
  r\in S.
\]

For every such \(r\), one has
\[
  \pi\equiv-1\pmod r.
\]
Since
\[
  r>5,
\]
in particular \(r\ne2\), and therefore
\[
  \ord_r(\pi)=2.
\]
\end{proof}

\begin{corollary}[Certified second-level reduction in the \(q=5\) branch]
\label{cor:q5-second-level-forced-reduction}
In the branch
\[
  \min S=5,
\]
the three child branches from the first \(5\)-input split admit the following
second-level certified reductions.

\begin{enumerate}[label=(\roman*),leftmargin=2.2em]

\item In \(\mathbf{Q5N}\), with non-Euler witness \(p\), the split
      \[
        3\mid C_5(p)
        \qquad\text{or}\qquad
        3\nmid C_5(p)
      \]
      is complete.  The first child is refuted by lower-prime avoidance.  The
      second child forces a support prime
      \[
        r>5,
        \qquad
        \ord_r(p)=5,
        \qquad
        r\equiv1\pmod5.
      \]

\item In \(\mathbf{Q5E1}\), with Euler prime \(\pi\), the split
      \[
        3\mid C_5(\pi)
        \qquad\text{or}\qquad
        3\nmid C_5(\pi)
      \]
      is complete.  The first child is refuted by lower-prime avoidance.  The
      second child forces a support prime
      \[
        r>5,
        \qquad
        \ord_r(\pi)=5,
        \qquad
        r\equiv1\pmod5.
      \]

\item In \(\mathbf{Q5E2}\), the split
      \[
        3\mid C_2(\pi)
        \qquad\text{or}\qquad
        3\nmid C_2(\pi)
      \]
      is complete.  The first child is refuted by lower-prime avoidance.  The
      second child forces a support prime
      \[
        r>5,
        \qquad
        \ord_r(\pi)=2.
      \]
\end{enumerate}
\end{corollary}

\begin{proof}
The three cases are exactly Propositions~\ref{prop:q5-phi5-second-level-split}
and \ref{prop:q5e2-second-level-split}, applied to the corresponding child
branches.
\end{proof}

\begin{definition}[Exponent schemas for second-level forced primes in the \(q=5\) branch]
\label{def:q5-second-level-forced-prime-schemas}
Let
\[
  r>5
\]
be a prime forced by Corollary~\ref{cor:q5-second-level-forced-reduction}.

\begin{enumerate}[label=(\roman*),leftmargin=2.2em]

\item If \(r\) is non-Euler, then its exponent satisfies
      \[
        e_r\equiv0\pmod2.
      \]
      Since \(\min S=5\), the lower-prime avoidance condition for \(3\) must
      also hold:
      \[
        \operatorname{Avoid}_3(r,e_r).
      \]
      Equivalently,
      \[
        \begin{cases}
          3\nmid e_r+1, & r\equiv1\pmod3,\\[4pt]
          \text{automatic}, & r\equiv2\pmod3.
        \end{cases}
      \]

\item If \(r\) is the Euler prime, then
      \[
        r\equiv1\pmod4,
        \qquad
        e_r=\alpha\equiv1\pmod4.
      \]
      In the branch \(\min S=5\), Proposition~\ref{prop:q5-euler-2mod3-empty}
      also forces
      \[
        r\equiv1\pmod3,
      \]
      and lower-prime avoidance gives
      \[
        3\nmid \alpha+1.
      \]
\end{enumerate}

In the branches \(\mathbf{Q5E1}\) and \(\mathbf{Q5E2}\), the forced prime
\(r\) is automatically non-Euler, because the Euler prime is already \(\pi\).
In the branch \(\mathbf{Q5N}\), the forced prime \(r\) may be either non-Euler
or the Euler prime, and this must be recorded as a further role split.
\end{definition}

\begin{corollary}[Role split for the forced prime in \(\mathbf{Q5N}\)]
\label{cor:q5n-forced-prime-role-split}
In the child branch \(\mathbf{Q5N}\) with
\[
  3\nmid C_5(p),
\]
let
\[
  r>5,
  \qquad
  r\mid C_5(p),
  \qquad
  \ord_r(p)=5
\]
be a forced support prime.  Then the split
\[
  r=\pi
  \qquad\text{or}\qquad
  r\ne\pi
\]
is complete.

If
\[
  r=\pi,
\]
then
\[
  r\equiv1\pmod4,
  \qquad
  \alpha\equiv1\pmod4,
  \qquad
  r\equiv1\pmod3,
  \qquad
  3\nmid\alpha+1.
\]
If
\[
  r\ne\pi,
\]
then \(r\) is non-Euler and
\[
  e_r\equiv0\pmod2,
  \qquad
  \operatorname{Avoid}_3(r,e_r)
\]
must hold.
\end{corollary}

\begin{proof}
Every support prime is either the Euler prime or non-Euler, and Euler's form
allows exactly one Euler prime.  Thus the split
\[
  r=\pi
  \qquad\text{or}\qquad
  r\ne\pi
\]
is exhaustive.

If \(r=\pi\), then Euler's form gives
\[
  r\equiv1\pmod4,
  \qquad
  \alpha\equiv1\pmod4.
\]
In the branch \(\min S=5\), Proposition~\ref{prop:q5-euler-2mod3-empty}
excludes the Euler residue
\[
  \pi\equiv2\pmod3.
\]
Hence
\[
  r=\pi\equiv1\pmod3.
\]
Since
\[
  \ord_3(\pi)=1,
\]
lower-prime avoidance gives
\[
  3\nmid\alpha+1.
\]

If \(r\ne\pi\), then \(r\) is non-Euler.  Euler's form gives
\[
  e_r\equiv0\pmod2.
\]
Since
\[
  \min S=5,
\]
the lower-prime avoidance condition
\[
  \operatorname{Avoid}_3(r,e_r)
\]
must hold.
\end{proof}

\begin{remark}[Certificate data for the second-level \(q=5\) splits]
A certificate for the second-level \(q=5\) split records one of the following
types:
\[
  \mathtt{Q5N\_Phi5\_second},
  \qquad
  \mathtt{Q5E1\_Phi5\_second},
  \qquad
  \mathtt{Q5E2\_Phi2\_second}.
\]

For a \(\Phi_5\)-record the verifier checks
\[
  C_5(x)=\frac{\Phi_5(x)}5,
  \qquad
  \Phi_5(x)=x^4+x^3+x^2+x+1,
\]
and verifies the split
\[
  3\mid C_5(x)
  \qquad\text{or}\qquad
  3\nmid C_5(x).
\]
If \(3\mid C_5(x)\), the certificate records a lower-prime avoidance
refutation.  If \(3\nmid C_5(x)\), the certificate records a prime divisor
\[
  r>5
\]
of \(C_5(x)\), verifies
\[
  \ord_r(x)=5,
\]
and records \(r\) as a forced support prime.

For a \(\Phi_2\)-record in \(\mathbf{Q5E2}\), the verifier checks
\[
  C_2(\pi)=\frac{\pi+1}{5^{v_5(\pi+1)}}
\]
and verifies the split
\[
  3\mid C_2(\pi)
  \qquad\text{or}\qquad
  3\nmid C_2(\pi).
\]
If \(3\mid C_2(\pi)\), the certificate records a lower-prime avoidance
refutation.  If \(3\nmid C_2(\pi)\), it records a prime divisor
\[
  r>5
\]
of \(C_2(\pi)\), verifies
\[
  \ord_r(\pi)=2,
\]
and records \(r\) as a forced support prime.

For forced primes in \(\mathbf{Q5N}\), the certificate must additionally record
the role split
\[
  r=\pi
  \qquad\text{or}\qquad
  r\ne\pi.
\]
For forced primes in \(\mathbf{Q5E1}\) and \(\mathbf{Q5E2}\), the forced prime
is automatically non-Euler.
\end{remark}


\subsection{Closure of the release \texorpdfstring{\(q=5\)}{q=5} branch}
\label{sec:q5-branch-closure}

We now prove the first complete small-prime branch closure in the present
certificate system.  Throughout this section we work under the minimal-prime
hypothesis \(q=\min\{p:p\mid N\}=5\).
The first \(5\)-input coverage split reduces this branch to the three child
inventories \(\mathbf{Q5E2}\), \(\mathbf{Q5E1}\), and \(\mathbf{Q5N}\).
The split is a coverage split, not a disjoint split: every hypothetical odd
perfect number with smallest prime divisor \(5\) realizes at least one of these
three children.

The purpose of this section is to close the three child inventories separately
and then combine them into a closure theorem for the \(\CertRelease\) \(q=5\)
branch inventory.  The child closures are proved in the order
\(\mathbf{Q5E2}\), then \(\mathbf{Q5E1}\), then \(\mathbf{Q5N}\).
Finally, Theorem~\ref{thm:q5-current-branch-inventory-closed} combines the three
closures with the first \(5\)-input coverage split.


\subsection{Closure of the \(\mathbf{Q5E2}\) successor inventory}

We first close the order-two Euler-input child branch \(\mathbf{Q5E2}\), where
\(\pi\equiv-1\pmod5\).
After the frozen first-window archive, the release terminal
\(\mathbf{Q5E2}\) successor inventory has four post-window pairs:
\((109,11)\), \((229,23)\), \((349,7)\), and \((409,41)\).
The finite part of each pair is checked by an aggregation-safe successor archive, while the infinite tail is closed by a parametric Zsigmondy--tail certificate.

\begin{table}[ht]
\centering
\small
\begin{tabular}{c|c|c|c}
\((\pi,r)\) & finite window & parametric window & status \\
\hline
\((109,11)\) &
\(e_{11}=22,\ldots,40\) &
all even \(e_{11}\ge22\) &
closed \\
\hline
\((229,23)\) &
\(e_{23}=22,\ldots,40\) &
all even \(e_{23}\ge22\) &
closed \\
\hline
\((349,7)\) &
\(e_7=22,\ldots,92\) &
all even \(e_7\ge94\) &
closed \\
\hline
\((409,41)\) &
\(e_{41}=22,\ldots,40\) &
all even \(e_{41}\ge22\) &
closed \\
\end{tabular}
\caption{Final \(\mathbf{Q5E2}\) post-window frontier.  All four
post-window pairs occurring in the \(\CertRelease\) \(\mathbf{Q5E2}\) successor
inventory are closed by finite aggregation-safe successor archives together
with parametric post-window certificates.}
\label{tab:q5e2-final-post-window-frontier}
\end{table}
\FloatBarrier

\begin{criterion}[Parametric post-window closure for \((\pi,r)=(349,7)\)]
\label{crit:q5e2-pi349-r7-parametric-post-window}
Work in the \(\mathbf{Q5E2}\) successor branch with \((\pi,r)=(349,7)\).
Let \(e_7\ge94\) be even and put \(n:=e_7+1\).  Then the corresponding
post-window row is closed.  If \(3\mid n\), lower-prime avoidance refutes the
row, since \(\operatorname{ord}_3(7)=1\) and hence
\(3\mid\sigma(7^{e_7})\).  If \(3\nmid n\), Zsigmondy's theorem applied to
\(7^n-1\) gives a primitive prime divisor \(s\mid7^n-1\) with
\(\operatorname{ord}_s(7)=n\).  Thus
\(s\mid\Phi_n(7)\mid\sigma(7^{e_7})\), so \(s\) is forced into the support.
The order relation \(n\mid s-1\), with \(n\) odd, gives
\(s\ge2n+1\).  Since \(e_7\ge94\), one has \(n\ge95\) and hence \(s\ge191\).
Therefore
\[
  H(\{5,7,349,s\})
  \left(\frac{59}{58}\right)^{18}<2.
\]
The row is consequently tail-controlled with \(B=59\) and \(M=18\).
\end{criterion}

\begin{proof}
The case \(3\mid n\) follows from lower-prime avoidance: \(7\equiv1\pmod3\),
so \(\operatorname{ord}_3(7)=1\), and \(3\mid n=e_7+1\) implies
\(3\mid\sigma(7^{e_7})\), contradicting \(3\notin S\).

Assume \(3\nmid n\).  Since \(e_7\) is even, \(n\) is odd, and \(n\ge95\); the
exceptional cases in Zsigmondy's theorem therefore do not apply to \(7^n-1\).
Let \(s\) be a primitive prime divisor.  Then
\(\operatorname{ord}_s(7)=n\) and \(s\mid\Phi_n(7)\mid\sigma(7^{e_7})\), so
\(s\) is forced into the support.  Moreover \(n\mid s-1\), and because \(n\)
and \(s\) are odd this gives \(s\ge2n+1\ge191\).
The verifier-facing tail inequality
\[
  H(\{5,7,349,s\})
  \left(\frac{59}{58}\right)^{18}<2
\]
is monotone decreasing in \(s\), and it holds already at \(s=191\).  Hence it
holds for every such primitive forced prime \(s\).  The row is tail-controlled.
\end{proof}

\begin{criterion}[Parametric post-window closure for \texorpdfstring{\((\pi,r)=(109,11)\)}{(pi,r)=(109,11)}]
\label{crit:q5e2-pi109-r11-parametric-post-window}
Work in the \(\mathbf{Q5E2}\) successor branch with \((\pi,r)=(109,11)\).
Let \(e_{11}\ge22\) be even and put \(n:=e_{11}+1\).  Then the corresponding
post-window row is tail-controlled.  Since \(e_{11}\) is even, \(n\) is odd,
and the divisor-sum factor is
\[
  \sigma(11^{e_{11}})
  =
  \frac{11^n-1}{10}
  =
  \prod_{\substack{d\mid n\\ d>1}}\Phi_d(11).
\]
Zsigmondy's theorem applied to \(11^n-1\) gives a primitive prime divisor
\(s\mid11^n-1\) with \(\operatorname{ord}_s(11)=n\), equivalently
\(s\mid\Phi_n(11)\).  Hence \(s\) is forced into the support.  The order
relation \(n\mid s-1\), together with the oddness of \(n\) and \(s\), gives
\(s\ge2n+1\).  Since \(e_{11}\ge22\), this yields \(n\ge23\) and \(s\ge47\).

The tail-control inequality
\[
  H(\{5,11,109,s\})
  \left(\frac{59}{58}\right)^{18}<2
\]
is monotone decreasing in \(s\), and it already holds for \(s=47\).  Therefore
it holds for every primitive forced prime \(s\) arising from \(\Phi_n(11)\).
Consequently every even post-window row \(e_{11}\ge22\) is tail-controlled with
\(B=59\) and \(M=18\).
\end{criterion}

\begin{proof}
The identity
\[
  \sigma(11^{e_{11}})
  =
  \frac{11^{e_{11}+1}-1}{11-1}
\]
gives the cyclotomic factorization after setting \(n=e_{11}+1\).

Since \(e_{11}\ge22\) and \(e_{11}\) is even, \(n\) is odd and at least \(23\).
Thus the exceptional cases in Zsigmondy's theorem do not apply.  Hence
\(11^n-1\) has a primitive prime divisor \(s\).  By primitivity,
\(\operatorname{ord}_s(11)=n\), so
\(s\mid\Phi_n(11)\mid\sigma(11^{e_{11}})\), and \(s\) is forced into the
support.  The relation \(n\mid s-1\), with \(n\) and \(s\) odd, gives
\(s\ge2n+1\ge47\).

Finally,
\[
  H(\{5,11,109,s\})
  =
  \frac54\cdot\frac{11}{10}\cdot\frac{109}{108}\cdot\frac{s}{s-1}.
\]
The expression is decreasing in \(s\).  Therefore it is enough to check the
endpoint \(s=47\), where the verifier confirms
\[
  H(\{5,11,109,47\})
  \left(\frac{59}{58}\right)^{18}<2.
\]
Hence all rows \(e_{11}\ge22\), \(e_{11}\equiv0\pmod2\), are tail-controlled.
\end{proof}

\begin{criterion}[Parametric post-window closure for \texorpdfstring{\((\pi,r)=(229,23)\)}{(pi,r)=(229,23)}]
\label{crit:q5e2-pi229-r23-parametric-post-window}
Work in the \(\mathbf{Q5E2}\) successor branch with \((\pi,r)=(229,23)\).
Let \(e_{23}\ge22\) be even and put \(n:=e_{23}+1\).  Then the corresponding
post-window row is tail-controlled.  Since \(e_{23}\) is even, \(n\) is odd,
and the divisor-sum factor is
\[
  \sigma(23^{e_{23}})
  =
  \frac{23^n-1}{23-1}
  =
  \frac{23^n-1}{22}.
\]
Equivalently, its odd cyclotomic part is contained in
\[
  \prod_{\substack{d\mid n\\ d>1}}\Phi_d(23).
\]
By Zsigmondy's theorem applied to \(23^n-1\), there is a primitive prime
divisor \(s\mid23^n-1\) with \(\operatorname{ord}_s(23)=n\), equivalently
\(s\mid\Phi_n(23)\).  Thus \(s\) divides \(\sigma(23^{e_{23}})\) and is forced
into the support.  Since \(n\mid s-1\), and since \(n\) and \(s\) are odd, one
has \(s\ge2n+1\).  From \(e_{23}\ge22\) we get \(n\ge23\), hence \(s\ge47\).

The verifier-facing tail-control inequality
\[
  H(\{5,23,229,s\})
  \left(\frac{59}{58}\right)^{18}<2
\]
is monotone decreasing in \(s\), and it already holds for \(s=19\).  Therefore
it certainly holds for every primitive forced prime \(s\ge47\) arising from
\(\Phi_n(23)\).  Consequently every even post-window row \(e_{23}\ge22\) is
tail-controlled with \(B=59\) and \(M=18\).
\end{criterion}

\begin{proof}
Set \(n=e_{23}+1\).  Since \(e_{23}\) is even and at least \(22\), \(n\) is odd
and \(n\ge23\).  The exceptional cases in Zsigmondy's theorem therefore do not
apply to \(23^n-1\).  Let \(s\) be a primitive prime divisor.  Then
\(\operatorname{ord}_s(23)=n\), so
\(s\mid\Phi_n(23)\mid\sigma(23^{e_{23}})\), and \(s\) is forced into the
support.  The relation \(n\mid s-1\), with \(n\) and \(s\) odd, gives
\(s\ge2n+1\ge47\).

Now
\[
  H(\{5,23,229,s\})
  =
  \frac54\cdot\frac{23}{22}\cdot\frac{229}{228}\cdot\frac{s}{s-1}.
\]
This expression is decreasing in \(s\).  Hence it is enough to check the
smallest permitted value.  In fact the stronger endpoint check
\[
  H(\{5,23,229,19\})
  \left(\frac{59}{58}\right)^{18}<2
\]
already holds.  Therefore the inequality also holds for every \(s\ge47\).
The row is tail-controlled with \(B=59\) and \(M=18\).
\end{proof}



\begin{criterion}[Parametric post-window closure for \texorpdfstring{\((\pi,r)=(409,41)\)}{(pi,r)=(409,41)}]
\label{crit:q5e2-pi409-r41-parametric-post-window}
Work in the \(\mathbf{Q5E2}\) successor branch with \((\pi,r)=(409,41)\).
Let \(e_{41}\ge22\) be even and put \(n:=e_{41}+1\).  Then the corresponding
post-window row is tail-controlled.  Since \(e_{41}\) is even, \(n\) is odd,
and the divisor-sum factor is
\[
  \sigma(41^{e_{41}})
  =
  \frac{41^n-1}{41-1}
  =
  \frac{41^n-1}{40}.
\]
Its primitive odd cyclotomic part is represented by
\[
  \Phi_n(41).
\]
By Zsigmondy's theorem applied to \(41^n-1\), there is a primitive prime
divisor \(s\mid41^n-1\) with \(\operatorname{ord}_s(41)=n\), equivalently
\(s\mid\Phi_n(41)\).  Thus \(s\) divides \(\sigma(41^{e_{41}})\) and is forced
into the support.  Since \(n\mid s-1\), and since \(n\) and \(s\) are odd, one
has \(s\ge2n+1\).  From \(e_{41}\ge22\) we get \(n\ge23\), hence \(s\ge47\).

The verifier-facing tail-control inequality
\[
  H(\{5,41,409,s\})
  \left(\frac{59}{58}\right)^{18}<2
\]
is monotone decreasing in \(s\), and it already holds for \(s=47\).  Therefore
it holds for every primitive forced prime \(s\ge47\) arising from
\(\Phi_n(41)\).  Consequently every even post-window row \(e_{41}\ge22\) is
tail-controlled with \(B=59\) and \(M=18\).
\end{criterion}

\begin{proof}
Set \(n=e_{41}+1\).  Since \(e_{41}\) is even and at least \(22\), \(n\) is odd
and \(n\ge23\).  The exceptional cases in Zsigmondy's theorem therefore do not
apply to \(41^n-1\).  Let \(s\) be a primitive prime divisor.  Then
\(\operatorname{ord}_s(41)=n\), so
\(s\mid\Phi_n(41)\mid\sigma(41^{e_{41}})\), and \(s\) is forced into the
support.  The relation \(n\mid s-1\), with \(n\) and \(s\) odd, gives
\(s\ge2n+1\ge47\).

Now
\[
  H(\{5,41,409,s\})
  =
  \frac54\cdot\frac{41}{40}\cdot\frac{409}{408}\cdot\frac{s}{s-1}.
\]
This expression is decreasing in \(s\).  Hence it is enough to check the
endpoint \(s=47\).  The verifier confirms
\[
  H(\{5,41,409,47\})
  \left(\frac{59}{58}\right)^{18}<2.
\]
Therefore the same inequality holds for every \(s\ge47\).  The row is
tail-controlled with \(B=59\) and \(M=18\).
\end{proof}


\begin{theorem}[Closure of the release \texorpdfstring{\(\mathbf{Q5E2}\)}{Q5E2} successor inventory]
\label{thm:q5e2-current-successor-inventory-closed}
In the \(\CertRelease\) \(\mathbf{Q5E2}\) successor inventory, the branch has no remaining
reduced, unresolved, or delegated terminal obligation.

More precisely, the terminal verified \(\mathbf{Q5E2}\) successor archive has
status
\[
  106/106\ \text{leaves checked},
\]
with status distribution
\[
  12\ \mathtt{lower\_prime\_refuted},
  \qquad
  16\ \mathtt{directly\_closed},
\]
\[
  10\ \mathtt{tail\_controlled},
  \qquad
  68\ \mathtt{child\_archive\_tail\_controlled}.
\]
It has no leaves with status
\[
  \mathtt{reduced}
  \qquad\text{or}\qquad
  \mathtt{unresolved}.
\]
Hence
\[
  \mathtt{recursive\_controlled}=\mathtt{true}.
\]

The post-window pairs occurring in the \(\CertRelease\) \(\mathbf{Q5E2}\) inventory are
exactly
\[
  (109,11),
  \qquad
  (229,23),
  \qquad
  (349,7),
  \qquad
  (409,41).
\]
Each of these pairs is closed by a finite aggregation-safe successor archive
together with a parametric post-window tail-control certificate.  Therefore the
\(\CertRelease\) \(\mathbf{Q5E2}\) successor inventory is exhausted.
\end{theorem}

\begin{proof}
The pair \((349,7)\) is closed by finite verification of the rows
\(22\le e_7\le92\), \(e_7\equiv0\pmod2\), together with the parametric
Zsigmondy--tail-control criterion for all even \(e_7\ge94\).

The pair \((109,11)\) is closed by the aggregation-safe \(E40\) successor
archive and by the parametric post-window criterion for all even
\(e_{11}\ge22\).  The corresponding primitive-prime argument uses
\(n=e_{11}+1\), \(s\mid\Phi_n(11)\), and \(\operatorname{ord}_s(11)=n\), so
that \(s\ge2n+1\ge47\).
The verifier checks
\[
  H(\{5,11,109,s\})
  \left(\frac{59}{58}\right)^{18}<2
\]
at the endpoint \(s=47\).

The pair \((229,23)\) is closed in the same way.  Its finite \(E40\) successor
archive is aggregation-safe, and its parametric certificate covers every even
\(e_{23}\ge22\).
The verifier checks the endpoint inequality
\[
  H(\{5,23,229,47\})
  \left(\frac{59}{58}\right)^{18}<2.
\]

The pair \((409,41)\) is also closed by an aggregation-safe \(E40\) successor
archive and a parametric certificate for all even \(e_{41}\ge22\).
The endpoint inequality checked by the verifier is
\[
  H(\{5,41,409,47\})
  \left(\frac{59}{58}\right)^{18}<2.
\]

The terminal archive contains no reduced or unresolved leaves, and the
machine-readable frontier inventory contains no additional \(\mathbf{Q5E2}\)
post-window pair outside the four pairs listed above.  Thus every terminal
obligation in the \(\CertRelease\) \(\mathbf{Q5E2}\) successor inventory is either
directly closed, lower-prime refuted, tail-controlled, or controlled by a
tail-controlled child archive.  This proves the claimed exhaustion of the
\(\CertRelease\) \(\mathbf{Q5E2}\) successor inventory.
\end{proof}

\begin{remark}[Scope of the \(\mathbf{Q5E2}\) exhaustion statement]
\label{rem:scope-q5e2-current-successor-inventory-closed}
Theorem~\ref{thm:q5e2-current-successor-inventory-closed} is an exhaustion
statement for the \(\CertRelease\) \(\mathbf{Q5E2}\) successor inventory.  It
closes the \(\mathbf{Q5E2}\) child branch, not by itself the parent
minimal-prime branch \(q=5\).

The parent \(q=5\) branch is closed only after combining this theorem with the
corresponding closures of \(\mathbf{Q5E1}\) and \(\mathbf{Q5N}\), and with the
certified first \(5\)-input coverage split.
\end{remark}

\subsection{Closure of the \texorpdfstring{\(\mathbf{Q5E1}\)}{Q5E1} successor inventory}
\label{subsec:q5e1-successor-inventory-closure}

We next close the order-one Euler-input child branch
\(\mathbf{Q5E1}\), where \(\pi\equiv1\pmod5\) and \(5\mid\alpha+1\).  In the
surviving \(q=5\) branch the Euler prime also satisfies
\(\pi\equiv1\pmod3\) and \(\pi\equiv1\pmod4\), hence every
\(\mathbf{Q5E1}\) Euler-prime candidate has \(\pi\equiv1\pmod{60}\).  The
finite exceptional window \(\pi<1381\), \(\pi\equiv1\pmod{60}\), is exhausted
by the frozen finite archive, while all remaining rows
\(\pi\ge1381\), \(\pi\equiv1\pmod{60}\), are closed by the parametric
\(C_5(\pi)\)-tail certificate below.



\begin{criterion}[Parametric Euler-prime window closure in \texorpdfstring{\(\mathbf{Q5E1}\)}{Q5E1}]
\label{crit:q5e1-parametric-pi-window-closure}
Work in the branch \(\mathbf{Q5E1}\).  Thus the Euler prime satisfies
\(\pi\equiv1\pmod5\) and \(5\mid\alpha+1\).  Together with the surviving
\(q=5\) constraints \(\pi\equiv1\pmod3\) and \(\pi\equiv1\pmod4\), this gives
\(\pi\equiv1\pmod{60}\).  Assume \(\pi\ge1381\).  Then the
\(\mathbf{Q5E1}\) row attached to \(\pi\) is tail-controlled.

More precisely, the forced cofactor is
\[
  C_5(\pi)
  :=
  \frac{\Phi_5(\pi)}5
  =
  \frac{\pi^4+\pi^3+\pi^2+\pi+1}{5}.
\]
Since
\(\pi\equiv1\pmod5\), one has \(v_5(\Phi_5(\pi))=1\), so \(C_5(\pi)\) is an
integer, and \(C_5(\pi)>1\).  Any prime divisor \(r\mid C_5(\pi)\) is forced
into the support.  The primes \(3\) and \(5\) do not divide
\(C_5(\pi)\), and every remaining prime divisor satisfies
\(r>5\), \(\operatorname{ord}_r(\pi)=5\), and \(r\equiv1\pmod5\); hence
\(r\ge11\).  The verifier checks the worst-case tail-control inequality at
\(\pi=1381\) and \(r=11\):
\[
  H(\{5,1381,11\})
  \left(\frac{59}{58}\right)^{18}
  <2.
\]
Since \(H(\{5,\pi,r\})\) is decreasing in both \(\pi\) and \(r\), this endpoint
check implies
\[
  H(\{5,\pi,r\})
  \left(\frac{59}{58}\right)^{18}
  <2
\]
for every \(\pi\ge1381\), \(\pi\equiv1\pmod{60}\), and
\(r\mid C_5(\pi)\).  Thus every \(\mathbf{Q5E1}\) row with \(\pi\ge1381\) is
tail-controlled with \(B=59\) and \(M=18\).
\end{criterion}

\begin{proof}
The congruence \(\pi\equiv1\pmod{60}\) follows from the \(\mathbf{Q5E1}\)
branch conditions together with the surviving \(q=5\) Euler-prime constraints.
Because \(5\mid\alpha+1\), the cyclotomic factor \(\Phi_5(\pi)\) divides
\(\sigma(\pi^\alpha)\).  Since \(\pi\equiv1\pmod5\), the standard LTE
calculation gives \(v_5(\Phi_5(\pi))=1\).  Thus
\(C_5(\pi)=\Phi_5(\pi)/5\) is an integer, greater than \(1\) for every
\(\pi\ge61\), hence in particular for every \(\pi\ge1381\).

Let \(r\mid C_5(\pi)\) be prime.  Then
\(r\mid\Phi_5(\pi)\mid\sigma(\pi^\alpha)\mid\sigma(N)=2N\).  Since \(r\) is
odd, \(r\mid N\), so \(r\) is forced into the support.

The full \(5\)-factor of \(\Phi_5(\pi)\) has been removed, so \(r\ne5\).
Also \(r\ne3\): because \(\pi\equiv1\pmod3\), one has
\[
  \Phi_5(\pi)
  =
  \pi^4+\pi^3+\pi^2+\pi+1
  \equiv
  5
  \equiv2
  \pmod3.
\]
Hence
\(3\nmid C_5(\pi)\).  Thus every prime divisor \(r\) of \(C_5(\pi)\) satisfies
\(r>5\).  For such an \(r\), since \(r\mid\Phi_5(\pi)\) and \(r\nmid5\), we
have \(\operatorname{ord}_r(\pi)=5\).  Therefore \(5\mid r-1\), so
\(r\equiv1\pmod5\), and the smallest possible such prime is \(r=11\).

Now
\[
  H(\{5,\pi,r\})
  =
  \frac54\cdot\frac{\pi}{\pi-1}\cdot\frac{r}{r-1}.
\]
This expression is decreasing in both \(\pi\) and \(r\).  Therefore the worst
case in the range \(\pi\ge1381\), \(r\ge11\), is \((\pi,r)=(1381,11)\).
The verifier checks
\[
  H(\{5,1381,11\})
  \left(\frac{59}{58}\right)^{18}
  =
  1.871757645581\ldots
  <2.
\]
Hence the same inequality holds throughout the whole parametric window.
\end{proof}

\begin{theorem}[Closure of the release \texorpdfstring{\(\mathbf{Q5E1}\)}{Q5E1} successor inventory]
\label{thm:q5e1-current-successor-inventory-closed}
The \(\CertRelease\) \(\mathbf{Q5E1}\) successor inventory is exhausted.

More precisely, the finite exceptional Euler-prime window
\[
  \pi<1381,
  \qquad
  \pi\equiv1\pmod{60},
\]
contains exactly the ten candidates
\[
  61,181,241,421,541,601,661,1021,1201,1321.
\]
These are precisely the leaves of the frozen first finite domain
\[
  \mathcal D_{5,E1}^{(1)}.
\]
That domain has verified status
\[
  10/10\ \text{leaves checked},
  \qquad
  10\ \mathtt{child\_archive\_tail\_controlled},
\]
with no reduced or unresolved leaves.  Hence
\[
  \mathtt{recursive\_controlled}=\mathtt{true}.
\]

For every remaining Euler-prime candidate
\[
  \pi\ge1381,
  \qquad
  \pi\equiv1\pmod{60},
\]
Criterion~\ref{crit:q5e1-parametric-pi-window-closure} gives a parametric
tail-control certificate with
\[
  B=59,
  \qquad
  M=18.
\]
Therefore the \(\CertRelease\) \(\mathbf{Q5E1}\) successor inventory has no remaining
reduced, unresolved, or delegated terminal obligation.
\end{theorem}

\begin{proof}
The finite window \(\pi<1381\), \(\pi\equiv1\pmod{60}\), is exhausted by direct
enumeration of the corresponding prime candidates:
\[
  61,181,241,421,541,601,661,1021,1201,1321.
\]
The frozen package \(\mathcal D_{5,E1}^{(1)}\) certifies all ten leaves as
\(\mathtt{child\_archive\_tail\_controlled}\).  For every remaining candidate
one has \(\pi\ge1381\) and \(\pi\equiv1\pmod{60}\).  The parametric criterion
applies and supplies a forced support prime \(r\mid C_5(\pi)\) with \(r\ge11\).
The endpoint inequality
\[
  H(\{5,1381,11\})
  \left(\frac{59}{58}\right)^{18}<2
\]
then implies the tail-control inequality for all larger \(\pi\) and all
admissible forced primes \(r\).  Hence every remaining row is tail-controlled.
Combining the finite exceptional window with the parametric window proves the
claimed exhaustion.
\end{proof}

\begin{remark}[Scope of the \(\mathbf{Q5E1}\) closure]
\label{rem:scope-q5e1-current-successor-inventory-closed}
Theorem~\ref{thm:q5e1-current-successor-inventory-closed} closes the
\(\CertRelease\) \(\mathbf{Q5E1}\) successor inventory.  It is one of the three child
closures needed for the \(q=5\) minimal-prime branch.

Together with the corresponding closures of \(\mathbf{Q5E2}\) and
\(\mathbf{Q5N}\), and with the first \(5\)-input coverage split, it contributes
to the final closure of the \(\CertRelease\) \(q=5\) branch inventory.
\end{remark}

\subsection{Closure of the \(\mathbf{Q5N}\) successor inventory}

We finally close the non-Euler \(5\)-input child branch
\(\mathbf{Q5N}\).  In this child branch there is a non-Euler witness prime
\(p\in S\), \(p\ne5\), \(p\ne\pi\), such that \(p\equiv1\pmod5\),
\(5\mid e_p+1\), and \(e_p\equiv0\pmod2\).  The finite exceptional witness
window \(p<211\), \(p\equiv1\pmod5\), is exhausted by the frozen finite archive,
while all remaining rows \(p\ge211\), \(p\equiv1\pmod5\), are closed by the
parametric \(C_5(p)\)-tail certificate below.


\begin{criterion}[Parametric witness-window closure in \texorpdfstring{\(\mathbf{Q5N}\)}{Q5N}]
\label{crit:q5n-parametric-witness-window-closure}
Work in the branch \(\mathbf{Q5N}\).  Thus there is a non-Euler \(5\)-input
witness \(p\in S\), \(p\ne5\), \(p\ne\pi\), with \(p\equiv1\pmod5\),
\(5\mid e_p+1\), and \(e_p\equiv0\pmod2\).  Assume \(p\ge211\).  Then the
\(\mathbf{Q5N}\) witness row attached to \(p\) is tail-controlled.

More precisely, the forced cofactor is
\[
  C_5(p)
  :=
  \frac{\Phi_5(p)}5
  =
  \frac{p^4+p^3+p^2+p+1}{5}.
\]
Since \(p\equiv1\pmod5\), one has \(v_5(\Phi_5(p))=1\), so \(C_5(p)\) is an
integer, and \(C_5(p)>1\).  Any prime divisor \(r\mid C_5(p)\) is forced into
the support.  The primes \(3\) and \(5\) do not divide
\(C_5(p)\), and every remaining prime divisor satisfies
\(r>5\), \(\operatorname{ord}_r(p)=5\), and \(r\equiv1\pmod5\); hence
\(r\ge11\).  The verifier checks the worst-case tail-control inequality at
\(p=211\) and \(r=11\):
\[
  H(\{5,211,11\})
  \left(\frac{59}{58}\right)^{18}
  <2.
\]
Since \(H(\{5,p,r\})\) is decreasing in both \(p\) and \(r\), this endpoint
check implies
\[
  H(\{5,p,r\})
  \left(\frac{59}{58}\right)^{18}
  <2
\]
for every \(p\ge211\), \(p\equiv1\pmod5\), and \(r\mid C_5(p)\).  Thus every
\(\mathbf{Q5N}\) witness row with \(p\ge211\) is tail-controlled with \(B=59\)
and \(M=18\).
\end{criterion}

\begin{proof}
Because \(5\mid e_p+1\), the cyclotomic factor \(\Phi_5(p)\) divides
\(\sigma(p^{e_p})\).  Since \(p\equiv1\pmod5\), the standard LTE calculation
gives \(v_5(\Phi_5(p))=1\).  Thus \(C_5(p)=\Phi_5(p)/5\) is an integer.

Let \(r\mid C_5(p)\) be prime.  Then
\(r\mid\Phi_5(p)\mid\sigma(p^{e_p})\mid\sigma(N)=2N\).  Since \(r\) is odd,
\(r\mid N\), so \(r\) is forced into the support.

The full \(5\)-factor of \(\Phi_5(p)\) has been removed, so
\(r\ne5\).  Also \(r\ne3\).  Indeed, in the \(\mathbf{Q5N}\) first witness window one has
\(p\equiv1\pmod5\), and the lower-prime avoidance condition for \(3\) rules out
the branch in which \(3\mid C_5(p)\).  Equivalently, every surviving
\(\mathbf{Q5N}\) witness row has \(3\nmid C_5(p)\).  Hence every surviving
prime divisor \(r\) of \(C_5(p)\) satisfies \(r>5\).  For such an \(r\), since
\(r\mid\Phi_5(p)\) and \(r\nmid5\), we have \(\operatorname{ord}_r(p)=5\).
Therefore \(5\mid r-1\), so \(r\equiv1\pmod5\), and the smallest possible such
prime is \(r=11\).

Now
\[
  H(\{5,p,r\})
  =
  \frac54\cdot\frac{p}{p-1}\cdot\frac{r}{r-1}.
\]
This expression is decreasing in both \(p\) and \(r\).  Therefore the worst case
in the range \(p\ge211\), \(r\ge11\), is \((p,r)=(211,11)\).
The verifier checks
\[
  H(\{5,211,11\})
  \left(\frac{59}{58}\right)^{18}
  =
  1.879308959140\ldots
  <2.
\]
Hence the same inequality holds throughout the whole parametric witness window.
\end{proof}

\begin{theorem}[Closure of the release \texorpdfstring{\(\mathbf{Q5N}\)}{Q5N} successor inventory]
\label{thm:q5n-current-successor-inventory-closed}
The \(\CertRelease\) \(\mathbf{Q5N}\) successor inventory is exhausted.

More precisely, the finite exceptional witness-prime window
\[
  p<211,
  \qquad
  p\equiv1\pmod5,
\]
contains exactly the ten witness candidates
\[
  11,31,41,61,71,101,131,151,181,191.
\]
These are precisely the leaves of the frozen first finite domain
\[
  \mathcal D_{5,N}^{(1)}.
\]
That domain has verified status
\[
  10/10\ \text{leaves checked},
  \qquad
  10\ \mathtt{child\_archive\_tail\_controlled},
\]
with no reduced or unresolved leaves.  Hence
\[
  \mathtt{recursive\_controlled}=\mathtt{true}.
\]

For every remaining witness-prime candidate
\[
  p\ge211,
  \qquad
  p\equiv1\pmod5,
\]
Criterion~\ref{crit:q5n-parametric-witness-window-closure} gives a parametric
tail-control certificate with
\[
  B=59,
  \qquad
  M=18.
\]
Therefore the \(\CertRelease\) \(\mathbf{Q5N}\) successor inventory has no remaining
reduced, unresolved, or delegated terminal obligation.
\end{theorem}

\begin{proof}
The finite window \(p<211\), \(p\equiv1\pmod5\), is exhausted by direct
enumeration of the corresponding prime candidates:
\[
  11,31,41,61,71,101,131,151,181,191.
\]
The frozen package \(\mathcal D_{5,N}^{(1)}\) certifies all ten leaves as
\(\mathtt{child\_archive\_tail\_controlled}\).  For every remaining candidate
one has \(p\ge211\) and \(p\equiv1\pmod5\).  The parametric criterion applies
and supplies a forced support prime \(r\mid C_5(p)\) with \(r\ge11\).
The endpoint inequality
\[
  H(\{5,211,11\})
  \left(\frac{59}{58}\right)^{18}<2
\]
then implies the tail-control inequality for all larger \(p\) and all admissible
forced primes \(r\).  Hence every remaining witness row is tail-controlled.

The possible Euler-role split \(r=\pi\) or \(r\ne\pi\) does not affect the
tail-control inequality, since the verified kernel \(\{5,p,r\}\) is already
sufficient in either role.  Combining the finite exceptional window with the
parametric witness window proves the claimed exhaustion.
\end{proof}

\begin{remark}[Scope of the \(\mathbf{Q5N}\) closure]
\label{rem:scope-q5n-current-successor-inventory-closed}
Theorem~\ref{thm:q5n-current-successor-inventory-closed} closes the
\(\CertRelease\) \(\mathbf{Q5N}\) successor inventory.  Together with the corresponding
closures of \(\mathbf{Q5E2}\) and \(\mathbf{Q5E1}\), and with the first
\(5\)-input coverage split, this completes the three child branches of the
\(\CertRelease\) \(q=5\) branch inventory.
\end{remark}

\subsection{Combination of the three child closures}
\label{subsec:q5-current-branch-inventory-closure}

\begin{theorem}[Closure of the release \texorpdfstring{\(q=5\)}{q=5} branch inventory]
\label{thm:q5-current-branch-inventory-closed}
The \(\CertRelease\) \(q=5\) branch inventory is exhausted.

Equivalently, within the certificate release \(\CertRelease\), there is no odd
perfect number \(N\) with
\[
  \min\{p:p\mid N\}=5.
\]
\end{theorem}

\begin{proof}
Assume \(\min\{p:p\mid N\}=5\).  The first necessary \(5\)-input coverage split
gives the three branch families \(\mathbf{Q5E2}\), \(\mathbf{Q5E1}\), and
\(\mathbf{Q5N}\).  This split is complete: every realization of the \(q=5\)
branch realizes at least one of these three children.

The \(\CertRelease\) \(\mathbf{Q5E2}\) successor inventory is exhausted by the
finite post-window closures for \((109,11)\), \((229,23)\), \((349,7)\), and
\((409,41)\), together with their parametric post-window certificates.

The \(\CertRelease\) \(\mathbf{Q5E1}\) successor inventory is exhausted by the
finite exceptional Euler-prime window \(\pi<1381\), \(\pi\equiv1\pmod{60}\),
together with the parametric Euler-prime window closure for \(\pi\ge1381\).

The \(\CertRelease\) \(\mathbf{Q5N}\) successor inventory is exhausted by the
finite exceptional witness-prime window \(p<211\), \(p\equiv1\pmod5\), together
with the parametric witness-prime window closure for \(p\ge211\).

Thus every child of the complete \(q=5\) first-input split is closed in the
\(\CertRelease\) successor inventory.  Hence the \(\CertRelease\) \(q=5\) branch
inventory is exhausted.
\end{proof}

\begin{remark}[Meaning of the \(q=5\) closure]
\label{rem:meaning-of-q5-closure}
Theorem~\ref{thm:q5-current-branch-inventory-closed} is a statement about the
minimal-prime branch \(q=\min\{p:p\mid N\}=5\).  It therefore excludes odd
perfect numbers whose smallest prime divisor is \(5\) inside the certificate
release \(\CertRelease\).

It should not be read as the statement that no odd perfect number is divisible
by \(5\); it excludes only the case in which \(5\) is the smallest prime
divisor.
\end{remark}




